\documentclass[conference]{IEEEtran}
\IEEEoverridecommandlockouts
\usepackage{cite}
\usepackage{amsmath,amssymb,amsfonts}
\usepackage{graphicx}
\usepackage{booktabs}
\usepackage{textcomp}
\usepackage{xcolor}
\usepackage{url}
\def\BibTeX{{\rm B\kern-.05em{\sc i\kern-.025em b}\kern-.08em
    T\kern-.1667em\lower.7ex\hbox{\sc e}\kern-.125em X}}

\begin{document}

\title{Attribution-First Measurement of GRPO Post-Training:\\
Signal Starvation, Stack Sensitivity, and a\\
Fourteen-Suite Held-Out Evaluation}

\author{\IEEEauthorblockN{Arvind C R}
\IEEEauthorblockA{\textit{Dept.\ of Computer Science and Engineering} \\
\textit{PES University}\\
Bengaluru, India \\
arvindcr4@gmail.com}
\and
\IEEEauthorblockN{Ramesh Prakash Guledgudd}
\IEEEauthorblockA{\textit{Dept.\ of Computer Science and Engineering} \\
\textit{PES University}\\
Bengaluru, India}
}

\maketitle

\begin{abstract}
Group Relative Policy Optimization (GRPO) has become a standard critic-free
method for reinforcement-learning post-training of large language models
(LLMs) with verifiable rewards. Yet ``GRPO'' names a family of
implementations whose parsers, samplers, regularisers and evaluators differ,
and its group-relative advantage silently discards every group whose
completions all receive the same reward. We present a multi-framework
benchmark harness and a measurement study built on one rule: every
comparative claim holds the whole training stack fixed except the factor
under test. We make three contributions. First, we formalise the
\emph{Zero-Variance Fraction} (ZVF), the share of a batch whose advantage is
identically zero, and show that it is the typical state of training
(ZVF $\approx$ 0.72--0.77 on our main configuration) and is mostly driven by
already-solved prompts. A signal model $S = p(1-p)(1-h_G(p))$ reproduces the
predicted U-shape across a 368-run audit. Second, we show that stack choice is
of the same order as algorithm choice: a variance decomposition attributes
$\eta^2 = 0.546$ to framework against $0.558$ to algorithm, and a matched
PPO/GRPO contrast is statistically indistinguishable (Welch $p = 0.7605$).
Third, we evaluate one GRPO-trained actor against fourteen held-out agentic
and reasoning suites under receipt-bound governance, obtaining, e.g.,
129/312 = 41.35\% pass@1 on VerilogEval and a reproduced 51.31\% on
Omni-MATH, while recording an explicit terminal state for every incomplete
lane. The results are deliberately conservative: no cross-suite aggregate is
computed and no improvement over a baseline is claimed.
\end{abstract}

\begin{IEEEkeywords}
reinforcement learning, GRPO, large language models, verifiable rewards,
reproducibility, benchmark evaluation
\end{IEEEkeywords}

\section{Introduction}

Reinforcement learning with verifiable rewards (RLVR) replaces the learned
reward model of RLHF~\cite{ouyang2022training} with a programmatic checker,
such as exact match for mathematics or unit tests for code. The workhorse
optimiser was long PPO~\cite{schulman2017proximal}, which needs a value
network. GRPO~\cite{shao2024deepseekmath,guo2025deepseekr1} removes the critic:
it samples a group of $G$ completions per prompt and standardises each reward
against the group's own statistics,
\begin{equation}
A_i = \frac{r_i - \mu_g}{\sigma_g + \varepsilon}, \quad
\mu_g = \tfrac{1}{G}\textstyle\sum_j r_j, \quad
\sigma_g^2 = \tfrac{1}{G}\textstyle\sum_j (r_j-\mu_g)^2 .
\label{eq:adv}
\end{equation}

This simplification has two consequences that standard training logs do not
reveal. The first concerns \emph{mechanism}. When every completion in a group
receives the same reward, $\sigma_g = 0$ and every $A_i$ in (\ref{eq:adv}) is
exactly zero. The group has been sampled and scored at full cost, but it
contributes no reward-driven gradient. The second concerns
\emph{attribution}. ``GRPO'' is implemented in Tinker~\cite{thinkingmachines2024tinker},
TRL~\cite{vonwerra2022trl}, veRL~\cite{sheng2024verl},
OpenRLHF~\cite{hu2024openrlhf} and many forks, each with its own answer
parser, temperature, KL coefficient, clip range, LoRA targets~\cite{hu2022lora}
and evaluator. Two runs labelled GRPO can differ in many ways that plausibly
change the outcome.

This paper reports an individual M.Tech project that treats both issues as
measurement problems. Our contributions are:
\begin{itemize}
\item a multi-framework harness with per-step telemetry and raw reward-tensor
persistence, so that every diagnostic is recomputable rather than read from a
log line;
\item the ZVF diagnostic, a falsifiable signal model for group size, and an
audit of an adaptive group-size controller derived from it;
\item evidence that the training stack is a material part of any reported
result, and a minimum reporting standard for stack-conditioned claims;
\item a governance-bound held-out evaluation of one trained actor on fourteen
suites, with every lane reported at its true terminal state.
\end{itemize}

\section{Related Work}

GRPO was introduced with DeepSeekMath~\cite{shao2024deepseekmath} and
popularised by DeepSeek-R1~\cite{guo2025deepseekr1}. Variants modify the
estimator: DAPO~\cite{yu2025dapo} adds decoupled clipping and dynamic sampling
that discards zero-variance groups, GSPO~\cite{qwen2025gspo} moves the ratio to
the sequence level, and Dr.\ GRPO~\cite{liu2025understanding} removes the
length and standard-deviation normalisations it identifies as biased.
Preference-based alternatives such as DPO~\cite{rafailov2024direct} and open
recipes such as T\"ulu~3~\cite{lambert2024tulu3} complete the landscape. On
methodology, Henderson et al.~\cite{henderson2018deep} showed that deep RL
results are highly sensitive to seeds and implementation details, and
reproducibility checklists~\cite{pineau2020improving}, model
cards~\cite{mitchell2019modelcards} and datasheets~\cite{gebru2021datasheets}
set norms for reporting. Our work differs in scope: we do not propose a new
estimator. We measure how much of the GRPO signal is actually used, and how
much of a reported difference belongs to the stack rather than the algorithm.

\section{Methodology}

\subsection{Attribution-first design}
Each comparison holds the model, prompts, verifier, decoding configuration,
step budget and evaluation fixed, and changes one factor. A shared task,
reward and decoding specification is dispatched through thin adapters to each
back-end, so the same experiment runs without re-implementing the grader.
Being able to run the same experiment on every back-end does not mean every
back-end completed it. In this release only Tinker and TRL produced completed
runs; veRL and OpenRLHF entries are dry-run launch plans and are never pooled
with real results.

\subsection{Telemetry and the Zero-Variance Fraction}
The training loop emits per-step ZVF, gradient utilisation, policy entropy,
completion length and KL divergence, and persists the raw per-group reward
tensors. For a batch of groups $\mathcal{B}$,
\begin{equation}
\mathrm{ZVF} = \frac{1}{|\mathcal{B}|}\sum_{g\in\mathcal{B}}
\mathbb{1}\left[\sigma_g = 0\right].
\end{equation}
Because it is computed from stored tensors, ZVF can be reproduced exactly.
We also track the pairwise contrast density (PCD), the fraction of
within-group completion pairs with differing rewards. PCD is less brittle
than ZVF under tiny reward perturbations.

\subsection{A signal model for group size}
For a prompt the policy solves with probability $p$, the probability that a
group of size $G$ is degenerate is $h_G(p) = p^G + (1-p)^G$. We model usable
signal per prompt as
\begin{equation}
S(p, G) = p(1-p)\,\bigl(1 - h_G(p)\bigr).
\label{eq:signal}
\end{equation}
This predicts (i) a U-shape of ZVF in accuracy, high at both extremes and
lowest at mid-range rewards; (ii) a monotone decrease in ZVF with $G$ at
interior accuracy; and (iii) a positive association between gradient norm and
$p(1-p)$. An adaptive controller uses (\ref{eq:signal}) to raise $G$ for
prompts whose groups keep collapsing.

\subsection{Statistical protocol}
We use multi-seed runs where cost allows, Welch and paired $t$-tests with
Benjamini--Hochberg correction across the test family, effect sizes with 95\%
confidence intervals, and power and minimum-detectable-effect (MDE)
accounting. We report held-out $n$ with every gain, and treat
inconclusive tests as inconclusive rather than as equivalence.

\subsection{Held-out evaluation campaign (E1--E14)}
One actor, \texttt{Qwen3.6-35B-A3B} with a Tinker-trained GRPO LoRA adapter
(bfloat16), is evaluated against fourteen suites. Each suite runs through its
own native evaluator at a pinned revision, with no per-suite tuning. Every
paid run carries a receipt (model, sampler path, revision, budget) and a
decontamination record. Where a suite's original contract is unavailable, a
declared replacement scope is run and reported separately. Original-contract
and replacement-scope numbers are never pooled.

\section{Results}

\subsection{Signal starvation is the typical regime}
On the main configuration ZVF is approximately 0.72--0.77, so roughly three
quarters of each rollout batch contributes no reward-driven gradient. The
collapsed groups are mostly \emph{all-correct}, i.e.\ prompts the model already
solves. Re-baselining the advantages does not recover useful signal; it only
turns the signal into a difficulty proxy. A curriculum that removes solved
prompts returned a null result on held-out accuracy. ZVF recomputed from three
stored GSM8K~\cite{cobbe2021training} tensors (0.130, 0.190, 0.155; pooled
0.158) matches the logged values. Its correlation with final held-out outcome
is weak ($\rho \approx 0.27$), so we present ZVF as a diagnostic, not a
predictor.

\subsection{Group size: a mechanism, not a universal optimum}
ZVF falls from 0.845 at $G=2$ to 0.631 at $G=16$, with diminishing returns
(Fig.~\ref{fig:gu}). A single-seed sweep does not establish an optimum. A
505-task conditional utility audit prefers $G=4$ over $G=5$ by only 0.00177
(95\% CI $[0.00082, 0.00270]$). The signal model explains why only a shallow
optimum should exist. Two of its three predictions reproduced: a 368-run audit
across seven models shows ZVF of 0.95--0.97 at both reward extremes and
0.25--0.29 at mid-range rewards (0.35--0.50), and a model$\times G$ grid shows
the monotone effect (Qwen3-32B: 0.33 $\to$ 0.18 $\to$ 0.08 across
$G = 4/8/16$). The gradient link, $\mathrm{corr}(\lVert\nabla\rVert,
p(1-p)) = +0.71$, was observed only on a 0.5B model and is directional.

\begin{figure*}[t]
\centering
\includegraphics[width=0.78\textwidth]{figures/fig_gradient_utilisation.pdf}
\caption{Gradient utilisation ($1-\mathrm{ZVF}$) rises with group size $G$
with diminishing returns (left); the relative gain per doubling of $G$
declines (right).}
\label{fig:gu}
\end{figure*}

\subsection{The adaptive controller}
In a four-arm open-trainer audit, adaptive group size matched the best
fixed-recipe held-out gain (+0.575) at 186 rollouts. DAPO-style dynamic
sampling drove ZVF to 0.00, but at +45\% rollout cost, which is an
accuracy-versus-compute trade rather than a controller win. In a
2,560-observation controller trace, 1,723 of 1,867 escalations (92.3\%) fired
on all-correct groups. A symmetric ZVF trigger therefore spends most of its
budget on already-solved prompts. We do not recommend the controller as a
default policy; the audit instead identifies the asymmetric trigger as the
next design target.

\subsection{The stack is part of the result}
Across 42 experiments, a variance decomposition attributes
$\eta^2 = 0.546$ of outcome variation to framework, $0.558$ to algorithm,
$0.471$ to model family and $0.322$ to scale. These factors are confounded,
but the framework term is clearly of the same order as the algorithm term. A
matched PPO/GRPO contrast on Qwen3-8B gives Welch $p = 0.7605$ (ten runs per
arm, Cohen's $d = -0.14$, 95\% CI $[-1.02, 0.74]$, power 0.06). This is a
failure to detect a difference, not evidence of equivalence. The same contrast on
Llama-3.1-8B-Instruct reverses to $d = 12.75$, and its arms ran on different
back-ends, so we withhold any directional algorithm claim. In a same-stack
audit of DAPO, GSPO, Dr.\ GRPO and AERO against GRPO (eight paired seeds, 500
held-out GSM8K items per unit), all four contrasts were inconclusive, and the
study's exact MDE (0.0101) exceeded its 0.010 design margin. Nominally matched
back-end runs reached final training rewards from 5.0\% to 85.6\%. Two runs
labelled ``DAPO'' showed mean ZVF of 0.00 and 0.58, because only one used
true dynamic sampling. Conversely, on one fixed stack four labelled
algorithms landed within a 0.034 band (0.710--0.744), which is comparable to
seed variation. When the stack is truly fixed, the labels stop separating.
These findings motivate an eight-item minimum reportable stack and a
machine-readable registry whose diff tool flags the ``DAPO'' pair as a
label-flip risk from manifests alone.

\subsection{Length bias}
Four of eleven GRPO runs peak before 65\% of training and then decay, which
is consistent with length bias rather than capability gain. Dr.\ GRPO does not
show this pattern, but its held-out gain is indistinguishable from GRPO's.
Mean completion length drifts from 194.4 to 183.5 tokens under a 200-token
cap, which limits what this comparison can show.

\subsection{Fourteen-suite held-out evaluation}
Table~\ref{tab:campaign} lists the lanes with graded results, and
Fig.~\ref{fig:lanes} shows the terminal state of all fourteen. Four scopes are
strictly complete: LAB-Bench~\cite{laurent2024labbench} public split
(1,967/1,967, micro accuracy 0.2288), AgentDojo~\cite{debenedetti2024agentdojo}
benign utility (97/97, utility 0.9072), VerilogEval~\cite{liu2023verilogeval}
(312/312, 41.35\% pass@1) and Omni-MATH~\cite{gao2024omnimath} (all 4,428
dispositions recorded, official accuracy 51.31\% reproduced). SWE-bench
Pro~\cite{deng2025swebenchpro} carries a full original-contract score of
2/731 = 0.274\% pass@1, with all 18 non-native outcomes kept in the
denominator. The remaining lanes are partial, quota-blocked or externally
blocked (e.g.\ private test sets), and each has a named reopen condition. A
separately authorised base-model run on BankerToolBench scored 0.0 over
100/100 trials. We trace this to tool-dialogue collapse (agents stopping after
2--14 steps, role-token degeneration), and report it as that rather than as a
measure of finance reasoning. Eleven of eleven deterministic ledger checks
(exact divisions, receipt presence, hash equality) passed.

\begin{table}[t]
\caption{Graded held-out results for one GRPO-trained actor. Each
suite uses its own native evaluator and denominator; values are not comparable
across rows and are never aggregated.}
\label{tab:campaign}
\centering
\small
\begin{tabular}{@{}llr@{}}
\toprule
Suite (scope) & Coverage & Score \\
\midrule
LAB-Bench (public) & 1,967/1,967 & 0.2288 \\
AgentDojo (benign utility) & 97/97 & 0.9072 \\
VerilogEval & 312/312 & 41.35\% \\
Omni-MATH & 4,428/4,428 & 51.31\% \\
SWE-bench Pro (original) & 731/731 & 0.274\% \\
\bottomrule
\end{tabular}
\end{table}

\begin{figure*}[t]
\centering
\includegraphics[width=0.82\textwidth]{figures/fig_lane_status.pdf}
\caption{Terminal state of every lane in the E1--E14 campaign: complete,
partial (with coverage), or blocked (quota or external), each with a named
reopen condition.}
\label{fig:lanes}
\end{figure*}

\section{Discussion and Threats to Validity}

The campaign results describe one actor at one revision; no per-suite
base-model comparator exists, so they are not evidence of what the GRPO step
contributed. Tinker runs are single-seed early-training snapshots of 30--50
steps, so asymptotic claims are unsupported, and Tinker results measure the
platform's GRPO implementation, whose internals are closed. Native
denominators differ across suites, which is why we report no average,
ranking or composite. The VerilogEval scores came through a managed sampler
path that no longer exists, so a re-run would change the sampler as well as
the run, even with the same weights. These limits are stated with the
numbers because hiding them would overstate the evidence.

\section{Conclusion}

At this compute scale the measurement layer is informative, while most
algorithmic levers are noise-limited. Signal starvation is the typical GRPO
regime, and group size has a mechanism and a shallow optimum rather than a
universal best value. The training stack moves results by as much as the
algorithm does, so a method label without its stack is not a reproducible
specification. The most durable outputs are tools that make results
checkable: the ZVF/PCD diagnostics, the minimum reportable stack and registry,
and a receipt-bound evaluation ledger. Future work will run a pre-registered
controller bake-off against a static $G = 16$ and a designed cross-framework
divergence study with identical task, reward, decoding and seeds across TRL,
veRL, OpenRLHF and Tinker.

\bibliographystyle{IEEEtran}
\bibliography{refs}

\end{document}
